%% slide12_study_note.tex
%% Detailed Study Note — Slide 12: Microgrid Mode-Switching Protection
%% TSA Meeting Preparation | Anoop V Eluvathingal | IIT Madras | 2026
\documentclass[12pt,a4paper]{article}

\usepackage[margin=2.5cm]{geometry}
\usepackage{amsmath,amssymb}
\usepackage{booktabs,tabularx,array}
\usepackage{graphicx}
\usepackage{xcolor}
\usepackage{tcolorbox}
\usepackage{enumitem}
\usepackage{fancyhdr}
\usepackage{titlesec}
\usepackage{hyperref}
\usepackage{parskip}

\definecolor{iitmnavy}{RGB}{15,40,80}
\definecolor{iitmgold}{RGB}{180,140,0}
\definecolor{lightblue}{RGB}{220,235,250}
\definecolor{lightyellow}{RGB}{255,252,220}
\definecolor{lightred}{RGB}{255,235,235}
\definecolor{lightgreen}{RGB}{230,255,230}

\tcbuselibrary{skins,breakable}
\newtcolorbox{keybox}[1]{
  colback=lightblue, colframe=iitmnavy, fonttitle=\bfseries\color{white},
  title=#1, arc=3pt, boxrule=1pt, breakable}
\newtcolorbox{formulabox}[1]{
  colback=lightyellow, colframe=iitmgold, fonttitle=\bfseries,
  title=#1, arc=3pt, boxrule=1pt, breakable}
\newtcolorbox{resultbox}[1]{
  colback=lightgreen, colframe=green!50!black, fonttitle=\bfseries,
  title=#1, arc=3pt, boxrule=1pt, breakable}
\newtcolorbox{warnbox}[1]{
  colback=lightred, colframe=red!60!black, fonttitle=\bfseries,
  title=#1, arc=3pt, boxrule=1pt, breakable}
\newtcolorbox{speakbox}{
  colback=orange!8, colframe=orange!70!black, fonttitle=\bfseries,
  title={What to Say at the TSA Meeting}, arc=3pt, boxrule=1pt, breakable}

\pagestyle{fancy}
\fancyhf{}
\fancyhead[L]{\color{iitmnavy}\small\textbf{TSA Meeting Study Note}}
\fancyhead[R]{\color{iitmnavy}\small Slide 12 — Microgrid Mode-Switching Protection}
\fancyfoot[C]{\thepage}
\fancyfoot[R]{\small\color{gray} Anoop V Eluvathingal $\cdot$ IIT Madras $\cdot$ 2026}
\renewcommand{\headrulewidth}{1.5pt}
\renewcommand{\headrule}{\color{iitmnavy}\hrule width\headwidth height\headrulewidth}

\titleformat{\section}{\large\bfseries\color{iitmnavy}}{}{0em}{}[\color{iitmnavy}\titlerule]
\titleformat{\subsection}{\normalsize\bfseries\color{iitmnavy!80}}{}{0em}{}

\hypersetup{colorlinks=true, linkcolor=iitmnavy, urlcolor=iitmnavy}

\begin{document}

%% ── Title Block ───────────────────────────────────────────────────────────────
\begin{tcolorbox}[colback=iitmnavy, coltext=white, arc=4pt, boxrule=0pt]
\begin{center}
{\Large\bfseries TSA Meeting Preparation — Detailed Study Note}\\[6pt]
{\large Slide 12: Ch 8 — Microgrid Mode-Switching Protection}\\[4pt]
{\normalsize SAMBP Framework $\cdot$ Anoop V Eluvathingal $\cdot$ IIT Madras $\cdot$ 2026}
\end{center}
\end{tcolorbox}

\vspace{0.5em}
\begin{center}
\small\textit{Sources: SAMBP TR-38, TR-39, TR-40, TR-41, TR-42 / 2026}
\end{center}

%% ── 1. Slide Overview ─────────────────────────────────────────────────────────
\section{1.\; Slide Overview and Purpose}

\begin{keybox}{What This Slide Is About}
This slide addresses one of the \textbf{most critical reliability gaps} in IBR
microgrid protection: when a microgrid transitions from \emph{grid-connected} to
\emph{islanded} mode, the fault current magnitude drops by \textbf{64--72\%} and
the current direction reference changes. Legacy relays designed for the grid-connected
state will either \textbf{mal-operate} or \textbf{fail to detect} faults in the
islanded state.

The slide presents the SAMBP solution: a \textbf{Bayesian mode detector} that
continuously classifies the operating state, and an \textbf{adaptive directional
element} that reconfigures itself in real time based on the detected mode.
\end{keybox}

\noindent\textbf{Position in the thesis:} Chapter 5 (Microgrid Protection), slide 12
of 22 in the TSA presentation. This is the second of two microgrid slides, following the
GFL/GFM fault current characterisation (slide 11).

%% ── 2. The Challenge ──────────────────────────────────────────────────────────
\section{2.\; The Engineering Challenge}

\subsection{Why the mode transition is dangerous}

When the Point-of-Common-Coupling (PCC) breaker opens (intentional or unintentional
islanding), the following changes occur simultaneously:

\begin{tabular}{@{}lp{5.5cm}p{5.5cm}@{}}
\toprule
\textbf{Parameter} & \textbf{Grid-Connected} & \textbf{Islanded}\\
\midrule
Fault current source & Upstream grid infeed ($\gg$IBR) & IBR only (limited to 1.1--2.0 pu)\\
Fault current magnitude & 14--22 pu (bolted) & 4--8 pu (IBR-limited)\\
Direction reference & Upstream grid voltage & GFM inverter terminal voltage\\
Frequency regulator & Stiff grid (50 Hz) & GFM virtual inertia ($H_{eq}=0.5$ s)\\
Distance relay validity & Valid (remote infeed) & \textbf{Invalid} (no remote reference)\\
87L validity & Valid (comms path open) & \textbf{Invalid} (PCC breaker severs link)\\
\bottomrule
\end{tabular}

\vspace{0.5em}
\begin{warnbox}{The 22\% Mal-Operation Problem}
In the SAMBP study, without mode-adaptive protection, \textbf{22\%} of fault events
during the grid-to-island transition window result in relay mal-operation (either
missed trip or false trip). This is the key ``before'' number that justifies the
entire microgrid protection chapter.
\end{warnbox}

%% ── 3. Mode Detector ──────────────────────────────────────────────────────────
\section{3.\; The Bayesian Mode Detector}

\subsection{Mathematical formulation}

The operating mode is classified at each sample using a Bayesian maximum likelihood
estimator:

\begin{formulabox}{Mode Detector Formula (slide equation)}
\begin{equation}
  \hat{m} = \arg\max_{m \in \{G, I\}} \Pr\!\bigl(z_t \mid \text{mode}=m,\; \hat{k}_\mathrm{ibr}\bigr)
  \label{eq:bayesian}
\end{equation}
where:
\begin{itemize}[nosep]
  \item $\hat{m}$ — estimated mode: $G$ = grid-connected, $I$ = islanded
  \item $z_t$ — measurement vector at time $t$ (frequency deviation, ROCOF,
        voltage magnitude, reactive-power perturbation response)
  \item $\hat{k}_\mathrm{ibr}$ — real-time IBR penetration estimate from the EKF
        digital twin (Chapter 6 / C3)
  \item $\Pr(\cdot)$ — conditional likelihood, updated at each 50 Hz sample
\end{itemize}
\end{formulabox}

\subsection{Why Bayesian, not a simple threshold?}

A fixed ROCOF threshold (IEEE 1547-2018: $\geq 1$ Hz/s) has a Non-Detection Zone
(NDZ) of $\pm 2\%$ rated power mismatch. The Bayesian estimator fuses four signals:

\begin{enumerate}[nosep]
  \item \textbf{ROCOF} — primary passive signal, $\geq 1$ Hz/s confirms islanding
        for $|\Delta P| \geq 0.02$ pu within 100 ms
  \item \textbf{Voltage/frequency window} — secondary check (V $\notin [0.88,1.10]$
        pu or $f \notin [49.0,50.5]$ Hz)
  \item \textbf{Active $\Delta Q$ injection} — 0.05 pu reactive pulse; grid-connected
        bus shows $\Delta f < 0.01$ Hz; islanded GFM shows $\Delta f = 0.10$ Hz
  \item \textbf{$\hat{k}_\mathrm{ibr}$ from EKF} — contextual prior (high $k_\mathrm{ibr}$
        means islanding is more likely given a given ROCOF value)
\end{enumerate}

The active Q-injection reduces the NDZ area by \textbf{87\%} (from $\pm 2\% \times
\pm 5\%$ to $\pm 1.2\% \times \pm 1.5\%$ in $\Delta P$--$\Delta Q$ space).

%% ── 4. Adaptive Directional Element ──────────────────────────────────────────
\section{4.\; Adaptive Directional Element (67)}

The slide lists two operating modes of the directional element:

\subsection{Grid-connected mode (G)}

\begin{itemize}[nosep]
  \item Polarising quantity: \textbf{pre-fault load flow voltage} from the upstream
        network — same as a conventional 67 relay
  \item Forward direction: toward the protected line (away from source)
  \item Reverse direction: toward the upstream grid
  \item Active elements: Z1/Z2 distance, 87L, 67Q/67N, Z2/POTT
  \item Pickup: $I_p = 1.50$ pu (Group 1 OC setting)
\end{itemize}

\subsection{Islanded mode (I)}

\begin{itemize}[nosep]
  \item Polarising quantity: \textbf{GFM inverter terminal voltage} (positive-sequence
        $V_1$ maintained by the GFM virtual impedance $\geq 0.5$ pu for all faults
        within the 1.1--1.5 pu current limit)
  \item Why: The upstream grid reference disappears when the PCC breaker opens;
        the GFM virtual impedance provides a stable local voltage reference
  \item Active elements: 87B (bus differential, always on), OC Group 2, 67 directional
  \item Disabled: Z1/Z2 distance, 87L, Z2/POTT (all require remote reference)
  \item Pickup: $I_p = 0.30$ pu (Group 2, 5$\times$ lower than grid-connected)
  \item \textbf{Key result:} OC Group 2 is 8$\times$ more sensitive than 87L
        ($k_\mathrm{ibr,min} = 0.0075$ pu vs 0.06 pu)
\end{itemize}

\subsection{The 20 ms mode-switch window}

Mode-switching is triggered by an \textbf{IEC 61850 GOOSE} signal from the TR-39
islanding relay, delivered within 4 ms of the islanding decision. The full sequence:

\begin{center}
\begin{tabular}{@{}lll@{}}
\toprule
\textbf{Event} & \textbf{Time} & \textbf{Protection state}\\
\midrule
PCC breaker opens (islanding) & $t = 0$ ms & Full grid-connected scheme\\
ROCOF reaches 1 Hz/s & $t = 100$ ms & ROCOF trip issued\\
GOOSE mode-switch signal & $t = 104$ ms & Within 4 ms of ROCOF trip\\
Setting group change complete & $t = 120$ ms & Islanded scheme active\\
\midrule
\textbf{Transition gap} & \textbf{20 ms} & 87B + 87L + OC-G1 still covering\\
\bottomrule
\end{tabular}
\end{center}

\textit{Zero protection gap}: during the 20 ms window, 87B (always active) and
87L (still enabled pre-switch) provide continuous fault coverage.

%% ── 5. Results ────────────────────────────────────────────────────────────────
\section{5.\; Results — What the Table in the Slide Means}

\begin{resultbox}{Mode-Switch Results Table Explained}
\begin{tabular}{@{}lp{3.5cm}p{3.5cm}p{4cm}@{}}
\toprule
\textbf{Metric} & \textbf{Fixed} & \textbf{SAMBP-MG} & \textbf{Explanation}\\
\midrule
Mode detection acc. & --- & \textbf{98.7\%} &
  Bayesian detector correctly identifies G vs I in 98.7\% of events
  (across all $\Delta P$, $\Delta Q$ combinations in the study)\\[4pt]
Transition mal. rate & 22\% & \textbf{1.4\%} &
  Relay mal-operations (false trip or missed trip) during the
  grid$\leftrightarrow$island transition; reduced by $15\times$\\[4pt]
Islanded fault TPR & 71\% & \textbf{97.3\%} &
  True positive rate for faults occurring in islanded mode;
  gap closed by OC Group 2 (lower pickup)\\[4pt]
False islands & 6/100 & \textbf{0/100} &
  Events where the scheme incorrectly declared islanding during
  normal operation; eliminated by the multi-signal Bayesian fusion\\
\bottomrule
\end{tabular}
\end{resultbox}

\subsection{The headline observation}

The slide states: \emph{``Mode-adaptive protection eliminates the 22\% mis-operation
rate seen at grid$\leftrightarrow$island transition boundary.''}

This is the \textbf{single most important result} of this slide. The 22\% figure
represents a very high risk for conventional relays and is the quantitative
justification for the entire adaptive microgrid protection framework.

%% ── 6. Supporting Figure ──────────────────────────────────────────────────────
\section{6.\; Best Supporting Figure — TR-39 Figure 3}

\begin{keybox}{Recommended Figure for the Presentation}
\textbf{Source:} SAMBP TR-39/2026, Figure 3 (page 3)\\
\textbf{Title:} Protection mode switching state machine\\[4pt]
The figure shows a state machine with five states:
\begin{enumerate}[nosep]
  \item \textbf{Grid-Connected Mode} — full scheme: Z1/Z2 + 87L + 67Q active, OC at grid level
  \item \textbf{Islanding Detection} — ROCOF/V/f monitoring active
  \item \textbf{Islanded Mode} — Z1/Z2 disabled, 87L offline, OC rescaled, GFM: V/f regulation
  \item \textbf{Resyncing (Relay 25)} — $\Delta V \leq 0.05$, $\Delta f \leq 0.10$ Hz,
        $\Delta\theta \leq 10^\circ$
  \item \textbf{Grid-Connected (restored)} or \textbf{Lockout} (resync failed)
\end{enumerate}
Timing annotations: detection $t_{det} = 100$ ms, mode switch $t_{sw} = 120$ ms from islanding,
Relay 25 window $\leq 5$ s.

\textbf{Why it is the best match:} It directly visualises the mode-transition logic
and timing that underlies every result in the slide table. It is self-contained and
immediately comprehensible to a TSA committee audience.
\end{keybox}

\textbf{Secondary figure (TR-40, Figure 3):} The three-state adaptive OC setting
group state machine (Group 1 GC $\to$ Group 2 Islanded $\to$ Group 3 Resync)
is an excellent companion to explain the adaptive directional element bullet points.

%% ── 7. Connection to Other Chapters ──────────────────────────────────────────
\section{7.\; Connections Across the Thesis}

\begin{tabular}{@{}lp{10cm}@{}}
\toprule
\textbf{Chapter/Section} & \textbf{Connection to this slide}\\
\midrule
Ch 2 (IBR Foundations) & GFL/GFM fault current model justifies the
  64--72\% reduction figure\\
Ch 3 (Differential) & 87L is \emph{disabled} in islanded mode (this slide);
  87B replaces it\\
Ch 4 (OC/Distance) & Distance relay Z1/Z2 also \emph{disabled} in
  islanded mode; OC Group 2 takes primary role\\
Ch 6 (EKF Digital Twin) & $\hat{k}_\mathrm{ibr}$ estimate feeds directly into the
  Bayesian mode detector (Eq.~\ref{eq:bayesian})\\
Ch 7 (WAPC) & IEC 61850 GOOSE is the triggering mechanism for the
  setting group change (4 ms delivery)\\
\bottomrule
\end{tabular}

%% ── 8. Equations Quick Reference ─────────────────────────────────────────────
\section{8.\; Equations Quick Reference}

\begin{formulabox}{Key Equations for This Slide}
\textbf{ROCOF-based islanding detection (TR-39):}
\begin{equation}
  \text{ROCOF} = \frac{|\Delta P| \cdot f_0}{2 H_{eq}} \geq 1 \;\text{Hz/s}
  \quad \Rightarrow \quad |\Delta P| \geq 0.02 \;\text{pu}
\end{equation}

\textbf{Active Q-injection frequency response (TR-39):}
\begin{equation}
  \Delta f = K_{f,Q} \cdot \Delta Q = 2.0 \times 0.05 = 0.10 \;\text{Hz}
  \quad (\text{detectable vs. floor } \Delta f_{resol} = 0.01 \;\text{Hz})
\end{equation}

\textbf{Islanded OC minimum sensitivity (TR-40):}
\begin{equation}
  k_{\min,\text{OC}} = \frac{I_{p,2} \times Z_\text{line}}{2}
  = \frac{0.30 \times 0.050}{2} = 0.0075 \;\text{pu}
  \quad (\text{8$\times$ more sensitive than 87L})
\end{equation}

\textbf{Resync transient energy check (TR-39):}
\begin{equation}
  \Delta W = \tfrac{1}{2} S_\text{ibr}(1 - \cos\Delta\theta_\text{max})
  = 0.0076 \;\text{pu$\cdot$s} \ll \Delta W_\text{limit} \approx 0.05 \;\text{pu$\cdot$s}
\end{equation}
\end{formulabox}

%% ── 9. What to Say ───────────────────────────────────────────────────────────
\section{9.\; TSA Meeting Speaker Script}

\begin{speakbox}
``This slide addresses what I call the \textbf{mode-transition blind spot} —
the period when the microgrid is switching from grid-connected to islanded operation.
During this window, legacy relays face a dilemma: they were set for high grid-fed
fault currents, but now the fault current has dropped by 64 to 72 percent, and the
directional reference has moved from the upstream grid to the GFM inverter terminal.

In our study, without adaptation, \textbf{22 percent} of events at the transition
boundary resulted in relay mal-operation. That is an unacceptably high failure rate.

Our solution has two parts. First, a \textbf{Bayesian mode detector} that fuses four
signals — ROCOF, voltage-frequency window, an active reactive power perturbation,
and the IBR penetration estimate from the digital twin — to classify the operating
mode at every sample. Second, an \textbf{adaptive directional element} that
automatically switches from pre-fault load-flow polarisation to GFM terminal
voltage polarisation when islanding is confirmed.

The result: the mal-operation rate falls from 22 percent to \textbf{1.4 percent}.
Islanded fault detection improves from 71 to \textbf{97.3 percent}. And false
island declarations — which would cause unnecessary outages — drop from 6 to
\textbf{zero} in a hundred test events.

The switching is triggered by an IEC 61850 GOOSE signal, delivered within
4 milliseconds of the islanding decision, with a total transition time of 20
milliseconds. During that window, the bus differential element provides
uninterrupted protection coverage — so there is \textbf{zero protection gap}.''
\end{speakbox}

%% ── 10. Anticipated Committee Questions ──────────────────────────────────────
\section{10.\; Anticipated TSA Committee Questions and Answers}

\begin{enumerate}
\item \textbf{Q: How is the Bayesian mode detector different from a simple ROCOF
      relay?}\\
      A: A simple ROCOF relay has a fixed NDZ of $\pm 2\%$ rated power. Our Bayesian
      detector fuses ROCOF with active Q-injection and the EKF's $\hat{k}_\mathrm{ibr}$
      estimate, reducing the NDZ area by 87\% and eliminating all false declarations.

\item \textbf{Q: What happens if the GOOSE signal is delayed or lost?}\\
      A: The 87B bus differential is always active in both modes and requires no GOOSE
      signal. It covers busbar faults during the transition. For the 20 ms gap, the
      pre-switch settings (OC Group 1 + 87L) remain active, providing backup coverage.

\item \textbf{Q: The 1.4\% residual mal-operation rate — what causes it?}\\
      A: Both residual misses are \emph{structural} — they occur when $k_\mathrm{ibr}
      < 0.0075$ pu, below the absolute detection floor of the islanded OC element.
      This corresponds to an inverter operating below 0.75\% of rated current, which
      is outside the scope of any protection scheme. The grid-code recommendation
      is $k_\mathrm{ibr} \geq 0.10$ pu per unit.

\item \textbf{Q: Is hardware-in-the-loop validation planned?}\\
      A: Yes. The identified next step is HIL testing on a commercial IED (SEL-411L)
      with a real-time PSCAD simulation via the DRTS interface.

\item \textbf{Q: How does this interact with the WAPC in Chapter 7?}\\
      A: The WAPC receives the islanding status via IEC 61850 MMS and updates the
      wide-area settings accordingly. The GOOSE mode-switch signal originates at
      the local islanding relay (TR-39) and is also broadcast to the WAPC for
      N-1 contingency re-assessment.
\end{enumerate}

%% ── Footer note ───────────────────────────────────────────────────────────────
\vspace{2em}
\begin{center}
\small\color{gray}
Study note prepared from SAMBP Technical Reports TR-38 through TR-42 (2026).\\
Best supporting figure: TR-39/2026, Figure 3 — Protection Mode Switching State Machine.\\
File: \texttt{slide12\_study\_note.pdf} $\cdot$ Generated \today
\end{center}

\end{document}
